\documentclass[12pt,a4paper]{article}
\usepackage[utf8]{inputenc}
\usepackage[T1]{fontenc}
\usepackage[spanish]{babel}
\usepackage{geometry}
\usepackage{booktabs}
\usepackage{longtable}
\usepackage{array}
\usepackage{xcolor}
\usepackage{listings}
\usepackage{hyperref}
\usepackage{fancyhdr}
\usepackage{graphicx}
\usepackage{titlesec}
\usepackage{enumitem}
\usepackage{mdframed}

\geometry{margin=2.5cm}

\definecolor{primaryblue}{RGB}{33,150,243}
\definecolor{criticalred}{RGB}{244,67,54}
\definecolor{normalgreen}{RGB}{76,175,80}
\definecolor{warningorange}{RGB}{255,152,0}
\definecolor{codegray}{RGB}{245,245,245}
\definecolor{darkgray}{RGB}{55,55,55}
\definecolor{accentpurple}{RGB}{103,58,183}

\lstset{
  backgroundcolor=\color{codegray},
  basicstyle=\ttfamily\small,
  breaklines=true,
  frame=single,
  rulecolor=\color{gray},
  keywordstyle=\color{primaryblue}\bfseries,
  commentstyle=\color{normalgreen},
  stringstyle=\color{criticalred},
  showstringspaces=false,
  tabsize=2,
  captionpos=b,
}

\pagestyle{fancy}
\fancyhf{}
\rhead{EdgeLeak AI -- Actualización \#6}
\lhead{Documentación Técnica}
\cfoot{\thepage}

\hypersetup{
  colorlinks=true,
  linkcolor=primaryblue,
  urlcolor=primaryblue,
}

\titleformat{\section}{\Large\bfseries\color{primaryblue}}{}{0em}{\thesection.\enspace}
\titleformat{\subsection}{\large\bfseries\color{darkgray}}{}{0em}{\thesubsection.\enspace}
\titleformat{\subsubsection}{\normalsize\bfseries}{}{0em}{\thesubsubsection.\enspace}

\title{
  \Huge\textbf{Documentación Técnica}\\[0.5em]
  \Large\textcolor{primaryblue}{EdgeLeak AI}\\[0.3em]
  \large Actualización \#6: API Local HTTP, Sensor Fusion\\
  e Integración con ESP32
}
\author{Proyecto EdgeLeak AI}
\date{\today}

\begin{document}

\maketitle
\thispagestyle{fancy}

\begin{mdframed}[linecolor=primaryblue,linewidth=1.5pt]
\textbf{Resumen Ejecutivo:} Esta actualización introduce la capa de comunicación en tiempo
real entre el hardware (ESP32) y la aplicación Flutter. Se implementa un servidor HTTP
local (\texttt{LocalApiService}) que expone dos endpoints REST en el puerto 8080:
(1)~\texttt{POST /api/sensor} para la ingesta de lecturas de los sensores de ruido y flujo
desde el ESP32; (2)~\texttt{GET /api/history} para la exportación segura del historial
de fugas detectadas a proyectos externos, protegido mediante una API Key en el header
\texttt{x-api-key}. Adicionalmente, el \texttt{DashboardController} implementa la lógica
de \textbf{Sensor Fusion}, un buffer circular de 10 lecturas y un mecanismo de
\textit{rate-limiting} para controlar las llamadas a la API de Groq.
\end{mdframed}

\vspace{1em}
\tableofcontents
\newpage

% ============================================================
\section{Contexto y Motivación}
% ============================================================

Hasta la Actualización~\#5, el \texttt{DashboardController} operaba únicamente con datos
simulados generados internamente por un \texttt{Timer} periódico. La integración con el
hardware real del ESP32 requería un canal de comunicación bidireccional que:

\begin{itemize}[leftmargin=2em]
  \item Permitiese al ESP32 \textbf{entregar} sus lecturas de sensor a la aplicación.
  \item Permitiese a \textbf{proyectos externos} consumir el historial de alertas de forma
        segura y estructurada.
  \item Implementase la \textbf{lógica de Sensor Fusion} en la capa de la aplicación,
        combinando la señal de ruido (ADC) y la señal de flujo (sensor de efecto Hall).
  \item Controlase la tasa de llamadas a la API de Groq para respetar los límites de la
        cuenta gratuita.
\end{itemize}

\begin{longtable}{>{\bfseries}p{3cm} p{4cm} p{6.5cm}}
\caption{Necesidades que originaron la Actualización \#6 --- La tabla describe cada
  necesidad técnica identificada y la solución implementada que la resuelve.}
\label{tab:necesidades} \\
\toprule
\textbf{Necesidad} & \textbf{Origen} & \textbf{Solución implementada} \\
\midrule
\endfirsthead
\toprule
\textbf{Necesidad} & \textbf{Origen} & \textbf{Solución implementada} \\
\midrule
\endhead
\midrule
\multicolumn{3}{r}{\small Continúa en la siguiente página\ldots} \\
\endfoot
\bottomrule
\endlastfoot
Canal de ingesta de datos del ESP32 &
  El ESP32 no tiene forma de enviar lecturas a la app sin un servidor local &
  \texttt{POST /api/sensor} en \texttt{LocalApiService} \\[0.5em]
Exportación segura del historial &
  Proyecto universitario asociado necesita consumir los datos detectados &
  \texttt{GET /api/history} protegido por \texttt{x-api-key} \\[0.5em]
Lógica de Sensor Fusion &
  Un solo sensor genera falsos positivos; se requiere contexto cruzado &
  \texttt{procesarLecturaSensor()} en \texttt{DashboardController} \\[0.5em]
Control de tasa de llamadas a Groq &
  La cuenta gratuita de Groq tiene límites de solicitudes por minuto &
  Contador de lecturas consecutivas + cooldown de 3 minutos \\
\end{longtable}

% ============================================================
\section{Arquitectura de la Solución}
% ============================================================

La solución se distribuye en dos archivos principales que colaboran mediante el patrón
de \emph{callbacks}:

\begin{longtable}{>{\ttfamily}p{6cm} p{2.5cm} p{5.5cm}}
\caption{Archivos principales de la Actualización \#6 --- La tabla lista cada archivo
  nuevo o modificado junto con su rol dentro de la arquitectura MVC de EdgeLeak AI.}
\label{tab:archivos-arch} \\
\toprule
\textbf{Archivo} & \textbf{Capa MVC} & \textbf{Responsabilidad} \\
\midrule
\endfirsthead
\toprule
\textbf{Archivo} & \textbf{Capa MVC} & \textbf{Responsabilidad} \\
\midrule
\endhead
\bottomrule
\endlastfoot
lib/data/services/local\_api\_service.dart & Servicio (Modelo) &
  Servidor HTTP local, routing de peticiones, seguridad de API Key. \\[0.5em]
lib/controllers/dashboard\_controller.dart & Controlador &
  Sensor Fusion, buffer circular, rate-limit de Groq, notificación reactiva a la UI. \\
\end{longtable}

\subsection{Diagrama de flujo de datos}

\begin{mdframed}[linecolor=primaryblue,linewidth=1pt,backgroundcolor=blue!3]
\begin{verbatim}
ESP32 (cada 5 s)
  │  POST /api/sensor
  │  {"ruido": 1800, "flujo": 0.04}
  ▼
LocalApiService  ──────────────────────────────────────────────────
  │  onProcesarLectura(ruido, flujo)  [callback]
  ▼
DashboardController.procesarLecturaSensor()
  │  1. Sensor Fusion  → estadoActual ("Posible Fuga")
  │  2. Buffer circular (últimas 10 lecturas)
  │  3. Contador de fugas consecutivas >= 3?
  │       └── SÍ y cooldown expirado? → llamarGroqPorFuga()
  │                                          └── DatabaseService.insertarAlerta()
  │  notifyListeners()  → UI se actualiza
  ▼
LocalApiService  ← {"status": "ok", "estado": "Posible Fuga", ...}
  │  HTTP 200 OK
  ▼
ESP32 recibe confirmación

Proyecto externo (cualquier cliente HTTP)
  │  GET /api/history
  │  Header: x-api-key: edgeleak-share-2024
  ▼
LocalApiService  → DatabaseService.obtenerHistorial()
  │  HTTP 200 OK
  ▼
[{"id":1,"veredicto":"Fuga Detectada","severidad":"Crítica",...}, ...]
\end{verbatim}
\end{mdframed}

% ============================================================
\section{LocalApiService --- Servidor HTTP Local}
% ============================================================

\subsection{Descripción general}

\texttt{LocalApiService} es un servicio de la capa de datos (\texttt{lib/data/services/})
que encapsula un \texttt{HttpServer} de la biblioteca estándar de Dart (\texttt{dart:io}).
Se inicia automáticamente cuando el \texttt{DashboardController} es instanciado y se
detiene cuando el controlador es destruido (\texttt{dispose()}).

\begin{lstlisting}[language=Dart, caption={Estructura de LocalApiService}]
class LocalApiService {
  final ProcesarLecturaCallback _onProcesarLectura;
  final EstadoActualCallback    _onEstadoActual;
  final DatabaseService         _dbService;
  HttpServer? _server;

  Future<void> iniciar() async { /* bind en 0.0.0.0:8080 */ }
  Future<void> detener() async { /* close(force: true) */ }
}
\end{lstlisting}

\subsection{Configuración de red}

\begin{longtable}{p{4cm} p{9.5cm}}
\caption{Parámetros de red del servidor HTTP --- La tabla describe cada parámetro
  de configuración del servidor y el valor utilizado en EdgeLeak AI.}
\label{tab:red} \\
\toprule
\textbf{Parámetro} & \textbf{Valor y descripción} \\
\midrule
\endfirsthead
\toprule
\textbf{Parámetro} & \textbf{Valor y descripción} \\
\midrule
\endhead
\bottomrule
\endlastfoot
Puerto & \texttt{8080}. Puerto TCP estándar para APIs de desarrollo. \\[0.5em]
Dirección de enlace & \texttt{InternetAddress.anyIPv4} (\texttt{0.0.0.0}). Acepta
  conexiones desde cualquier interfaz de red del dispositivo (WiFi, Ethernet, USB
  tethering). Esto es necesario para que el ESP32 pueda alcanzar la app desde la
  misma red local. \\[0.5em]
Tipo de contenido & \texttt{application/json; charset=utf-8}. Todas las respuestas
  llevan este \texttt{Content-Type} declarado en el header. \\[0.5em]
API Key & Leída de la variable de entorno \texttt{LOCAL\_API\_KEY} en el archivo
  \texttt{.env}. Si la variable no está definida, se usa el valor de referencia
  \texttt{edgeleak-share-2024}. \\
\end{longtable}

% -----------------------------------------------------------
\subsection{Endpoint 1: POST /api/sensor}
% -----------------------------------------------------------

\subsubsection{Descripción}

Recibe la lectura periódica (cada 5 segundos) del ESP32 con los valores de los dos
sensores y desencadena la lógica de Sensor Fusion en el \texttt{DashboardController}.

\subsubsection{Especificación de la petición}

\begin{longtable}{p{3cm} p{10.5cm}}
\caption{Especificación de POST /api/sensor --- La tabla define cada aspecto del contrato
  HTTP del endpoint de ingesta de datos del sensor.}
\label{tab:spec-post} \\
\toprule
\textbf{Atributo} & \textbf{Valor} \\
\midrule
\endfirsthead
\toprule
\textbf{Atributo} & \textbf{Valor} \\
\midrule
\endhead
\bottomrule
\endlastfoot
URL & \texttt{http://<IP\_DEL\_TELEFONO>:8080/api/sensor} \\[0.5em]
Método HTTP & \texttt{POST} \\[0.5em]
Content-Type & \texttt{application/json} \\[0.5em]
Autenticación & Ninguna (endpoint público en la red local) \\[0.5em]
Cuerpo (body) & JSON con dos campos obligatorios:
  \texttt{"ruido"} (entero, valor ADC 0--4095) y
  \texttt{"flujo"} (decimal, caudal en L/min) \\[0.5em]
Frecuencia esperada & Cada 5 segundos desde el ESP32 \\
\end{longtable}

\subsubsection{Ejemplo de cuerpo de la petición}

\begin{lstlisting}[language={}  , caption={Payload JSON enviado por el ESP32}]
{
  "ruido": 1800,
  "flujo": 0.04
}
\end{lstlisting}

\subsubsection{Respuestas posibles}

\begin{longtable}{p{2.5cm} p{4cm} p{7cm}}
\caption{Respuestas de POST /api/sensor --- La tabla describe cada código de respuesta
  HTTP posible, su significado y el ejemplo de cuerpo devuelto.}
\label{tab:resp-post} \\
\toprule
\textbf{HTTP Status} & \textbf{Significado} & \textbf{Ejemplo de cuerpo} \\
\midrule
\endfirsthead
\toprule
\textbf{HTTP Status} & \textbf{Significado} & \textbf{Ejemplo de cuerpo} \\
\midrule
\endhead
\bottomrule
\endlastfoot
\texttt{200 OK} &
  Lectura procesada correctamente. &
  \texttt{\{"status":"ok","estado":"Posible Fuga","ruido":1800,"flujo":0.04\}} \\[0.5em]
\texttt{400 Bad Request} &
  El JSON no contiene los campos obligatorios \texttt{ruido} y \texttt{flujo}. &
  \texttt{\{"error":"Campos requeridos: ruido (int) y flujo (double)"\}} \\[0.5em]
\texttt{400 Bad Request} &
  El cuerpo de la petición no es un JSON válido. &
  \texttt{\{"error":"JSON inválido"\}} \\[0.5em]
\texttt{500 Internal Server Error} &
  Error inesperado en el servidor. &
  \texttt{\{"error":"Error interno del servidor"\}} \\
\end{longtable}

\subsubsection{Código de ejemplo para el ESP32 (Arduino C++)}

El siguiente fragmento muestra cómo programar el ESP32 para enviar las lecturas de
sensor a la app Flutter. La dirección IP debe ser la del teléfono dentro de la red
WiFi local.

\begin{lstlisting}[language=C++, caption={Código ESP32: envío periódico al POST /api/sensor}]
#include <WiFi.h>
#include <HTTPClient.h>
#include <ArduinoJson.h>

const char* ssid     = "NOMBRE_DE_TU_RED_WIFI";
const char* password = "CONTRASENA_WIFI";
const char* serverUrl = "http://192.168.1.100:8080/api/sensor";
// ^^ Reemplazar con la IP del telefono Android/iOS en tu red local

// Pin del sensor de flujo YF-S201 (interrupcion)
const int FLOW_SENSOR_PIN = 14;
volatile int pulseCount = 0;

void IRAM_ATTR onFlowPulse() { pulseCount++; }

void setup() {
  Serial.begin(115200);
  WiFi.begin(ssid, password);
  while (WiFi.status() != WL_CONNECTED) {
    delay(500);
    Serial.print(".");
  }
  Serial.println("\nWiFi conectado. IP: " + WiFi.localIP().toString());
  attachInterrupt(digitalPinToInterrupt(FLOW_SENSOR_PIN),
                  onFlowPulse, RISING);
}

void loop() {
  delay(5000); // Enviar cada 5 segundos

  // 1. Leer sensor de ruido (ADC, pin 34)
  int ruido = analogRead(34);

  // 2. Calcular flujo: YF-S201 = 7.5 pulsos por litro por minuto
  float flujoLPM = (pulseCount / 7.5);
  pulseCount = 0;

  // 3. Construir el JSON
  StaticJsonDocument<64> doc;
  doc["ruido"] = ruido;
  doc["flujo"] = flujoLPM;
  String payload;
  serializeJson(doc, payload);

  // 4. Enviar la peticion HTTP POST
  if (WiFi.status() == WL_CONNECTED) {
    HTTPClient http;
    http.begin(serverUrl);
    http.addHeader("Content-Type", "application/json");
    int httpCode = http.POST(payload);

    if (httpCode == HTTP_CODE_OK) {
      Serial.println("OK: " + http.getString());
    } else {
      Serial.println("Error HTTP: " + String(httpCode));
    }
    http.end();
  }
}
\end{lstlisting}

\begin{mdframed}[linecolor=warningorange,linewidth=1.5pt,backgroundcolor=orange!3]
\textbf{Cómo obtener la IP del teléfono:} En Android, ir a
\textit{Ajustes} $\rightarrow$ \textit{Conexiones} $\rightarrow$ \textit{WiFi}
$\rightarrow$ \textit{Nombre de la red} $\rightarrow$ \textit{Avanzado}.
En iOS, ir a \textit{Ajustes} $\rightarrow$ \textit{WiFi} $\rightarrow$
\textit{(i) junto a la red conectada}. La app y el ESP32 deben estar conectados
a la \textbf{misma red WiFi local}.
\end{mdframed}

% -----------------------------------------------------------
\subsection{Endpoint 2: GET /api/history}
% -----------------------------------------------------------

\subsubsection{Descripción}

Devuelve el historial completo de alertas de fuga almacenadas en la base de datos
SQLite local de la app. Este endpoint está diseñado para ser consumido por proyectos
externos (p.\,ej., un sistema universitario de análisis de datos) y está protegido
mediante una API Key que debe incluirse en el header \texttt{x-api-key}.

\subsubsection{Especificación de la petición}

\begin{longtable}{p{3cm} p{10.5cm}}
\caption{Especificación de GET /api/history --- La tabla define el contrato HTTP
  completo del endpoint de exportación, incluyendo autenticación y campos de respuesta.}
\label{tab:spec-get} \\
\toprule
\textbf{Atributo} & \textbf{Valor} \\
\midrule
\endfirsthead
\toprule
\textbf{Atributo} & \textbf{Valor} \\
\midrule
\endhead
\bottomrule
\endlastfoot
URL & \texttt{http://<IP\_DEL\_TELEFONO>:8080/api/history} \\[0.5em]
Método HTTP & \texttt{GET} \\[0.5em]
Autenticación & Header \texttt{x-api-key} con el valor de la clave compartida \\[0.5em]
Cuerpo de la petición & Ninguno \\[0.5em]
Clave por defecto & \texttt{edgeleak-share-2024} (configurable en \texttt{.env} con
  la variable \texttt{LOCAL\_API\_KEY}) \\
\end{longtable}

\subsubsection{Autenticación: API Key}

El header \texttt{x-api-key} es \textbf{obligatorio}. Si no se incluye o si su valor
no coincide con la clave configurada, el servidor devuelve \texttt{401 Unauthorized}:

\begin{lstlisting}[language={}, caption={Respuesta cuando la API Key es incorrecta o falta}]
HTTP/1.1 401 Unauthorized
Content-Type: application/json

{"error": "No autorizado. Header x-api-key invalido o ausente."}
\end{lstlisting}

\subsubsection{Ejemplo de petición correcta}

\begin{lstlisting}[language={}, caption={Petición GET /api/history con autenticación correcta}]
GET /api/history HTTP/1.1
Host: 192.168.1.100:8080
x-api-key: edgeleak-share-2024
\end{lstlisting}

\subsubsection{Estructura de la respuesta exitosa (200 OK)}

La respuesta es un array JSON donde cada elemento representa una alerta de fuga
registrada en la base de datos local. Los elementos están ordenados del más reciente
al más antiguo.

\begin{lstlisting}[language={}, caption={Ejemplo de respuesta exitosa de GET /api/history}]
HTTP/1.1 200 OK
Content-Type: application/json

[
  {
    "id": 5,
    "veredicto": "Fuga Detectada",
    "severidad": "Critica",
    "mensaje": "Reemplaza el latiguillo de inmediato.",
    "timestamp": "2024-11-15T14:32:10.000Z"
  },
  {
    "id": 4,
    "veredicto": "Flujo Normal",
    "severidad": "Normal",
    "mensaje": "El caudal se encuentra dentro de parametros.",
    "timestamp": "2024-11-15T12:10:05.000Z"
  }
]
\end{lstlisting}

\subsubsection{Campos del objeto de respuesta}

\begin{longtable}{>{\ttfamily}p{3cm} p{2.5cm} p{8cm}}
\caption{Campos de cada objeto en la respuesta de GET /api/history --- La tabla describe
  el tipo de dato y el significado de cada campo devuelto por el endpoint de historial.}
\label{tab:campos-history} \\
\toprule
\textbf{Campo} & \textbf{Tipo} & \textbf{Descripción} \\
\midrule
\endfirsthead
\toprule
\textbf{Campo} & \textbf{Tipo} & \textbf{Descripción} \\
\midrule
\endhead
\bottomrule
\endlastfoot
id & integer &
  Identificador único autoincremental de la alerta en la base de datos SQLite. \\[0.5em]
veredicto & string &
  Clasificación de la IA. Valores posibles: \texttt{"Fuga Detectada"} o
  \texttt{"Flujo Normal"}. \\[0.5em]
severidad & string &
  Nivel de urgencia devuelto por la IA. Valores posibles: \texttt{"Crítica"},
  \texttt{"Advertencia"} o \texttt{"Normal"}. \\[0.5em]
mensaje & string &
  Recomendación en lenguaje natural generada por el LLM de Groq. \\[0.5em]
timestamp & string &
  Fecha y hora de la detección en formato ISO~8601
  (\texttt{YYYY-MM-DDTHH:mm:ss.sssZ}). \\
\end{longtable}

\subsubsection{Respuestas posibles}

\begin{longtable}{p{2.5cm} p{4cm} p{7cm}}
\caption{Códigos de respuesta de GET /api/history --- La tabla describe cada código
  HTTP posible y las condiciones que lo generan.}
\label{tab:resp-get} \\
\toprule
\textbf{HTTP Status} & \textbf{Significado} & \textbf{Condición} \\
\midrule
\endfirsthead
\toprule
\textbf{HTTP Status} & \textbf{Significado} & \textbf{Condición} \\
\midrule
\endhead
\bottomrule
\endlastfoot
\texttt{200 OK} &
  Historial devuelto correctamente. &
  Header \texttt{x-api-key} presente y correcto. Si no hay alertas aún, devuelve
  un array vacío \texttt{[]}. \\[0.5em]
\texttt{401 Unauthorized} &
  Acceso denegado. &
  El header \texttt{x-api-key} no está presente o su valor no coincide con la clave. \\[0.5em]
\texttt{500 Internal Server Error} &
  Error del servidor. &
  Fallo inesperado al consultar la base de datos. \\
\end{longtable}

% -----------------------------------------------------------
\subsection{Respuestas de error generales}
% -----------------------------------------------------------

\begin{longtable}{p{2.5cm} p{11cm}}
\caption{Respuestas de error generales del servidor --- La tabla lista los errores
  que aplican a cualquier endpoint, independientemente del método HTTP.}
\label{tab:errores-generales} \\
\toprule
\textbf{HTTP Status} & \textbf{Condición y cuerpo de respuesta} \\
\midrule
\endfirsthead
\toprule
\textbf{HTTP Status} & \textbf{Condición y cuerpo de respuesta} \\
\midrule
\endhead
\bottomrule
\endlastfoot
\texttt{404 Not Found} &
  Se solicitó una ruta que no existe. Ej.: \texttt{GET /api/status}.
  Cuerpo: \texttt{\{"error":"Endpoint no encontrado: GET /api/status"\}} \\[0.5em]
\texttt{500 Internal Server Error} &
  Excepción no controlada durante el procesamiento de la petición.
  Cuerpo: \texttt{\{"error":"Error interno del servidor"\}} \\
\end{longtable}

% ============================================================
\section{DashboardController --- Cambios Implementados}
% ============================================================

\subsection{Nuevas propiedades}

\begin{longtable}{>{\ttfamily}p{5cm} p{2.5cm} p{6cm}}
\caption{Nuevas propiedades añadidas a DashboardController --- La tabla describe cada
  campo nuevo y su propósito dentro de la lógica de fusión de sensores y control de tasa.}
\label{tab:props-dashboard} \\
\toprule
\textbf{Propiedad} & \textbf{Tipo} & \textbf{Descripción} \\
\midrule
\endfirsthead
\toprule
\textbf{Propiedad} & \textbf{Tipo} & \textbf{Descripción} \\
\midrule
\endhead
\bottomrule
\endlastfoot
\_estadoActual & String & Estado de fusión de la última lectura. Getter público:
  \texttt{estadoActual}. \\[0.5em]
\_bufferLecturas & List\textless{}Map\textgreater & Buffer circular de las últimas
  10 lecturas. Getter público: \texttt{bufferLecturas} (no modificable). \\[0.5em]
\_contadorPosibleFuga & int & Contador de lecturas consecutivas con estado
  ``Posible Fuga''. Se reinicia cuando el estado cambia o cuando se llama a Groq. \\[0.5em]
\_lastGroqEvaluation & DateTime? & Timestamp de la última llamada a Groq. Nulo si
  nunca se ha llamado. Controla el cooldown de 3 minutos. \\[0.5em]
\_umbralRuido & int & Umbral de ruido (ADC) para Sensor Fusion. Valor: \texttt{1500}. \\[0.5em]
\_umbralFlujoNormal & double & Umbral de flujo (L/min) para clasificar ``Uso Normal''.
  Valor: \texttt{0.5}. \\[0.5em]
\_umbralFlujoFuga & double & Umbral de flujo (L/min) por debajo del cual se clasifica
  ``Posible Fuga''. Valor: \texttt{0.1}. \\
\end{longtable}

% -----------------------------------------------------------
\subsection{Lógica de Sensor Fusion}
% -----------------------------------------------------------

El método \texttt{procesarLecturaSensor(int ruido, double flujo)} implementa las reglas
de clasificación cruzando los dos sensores:

\begin{longtable}{p{3.5cm} p{3.5cm} p{6.5cm}}
\caption{Tabla de decisión de Sensor Fusion --- La tabla define las tres condiciones de
  clasificación y las acciones que desencadena cada una dentro del sistema.}
\label{tab:fusion} \\
\toprule
\textbf{Condición de ruido} & \textbf{Condición de flujo} & \textbf{Estado resultante y acción} \\
\midrule
\endfirsthead
\toprule
\textbf{Condición de ruido} & \textbf{Condición de flujo} & \textbf{Estado resultante y acción} \\
\midrule
\endhead
\bottomrule
\endlastfoot
\texttt{ruido > 1500} & \texttt{flujo > 0.5} &
  \textbf{Uso Normal} --- La tubería transporta agua activamente (grifo abierto).
  Sin alerta. Reinicia el contador de fugas. \\[0.5em]
\texttt{ruido > 1500} & \texttt{flujo < 0.1} &
  \textbf{Posible Fuga} --- Ruido elevado con flujo casi nulo. Incrementa el
  contador de fugas consecutivas. Si llega a 3, activa Groq (si el cooldown
  ha expirado). \\[0.5em]
Cualquier otro caso & Cualquier otro caso &
  \textbf{Sin Clasificar} --- La combinación de señales no corresponde a ningún
  patrón definido. Reinicia el contador de fugas. \\
\end{longtable}

\begin{lstlisting}[language=Dart, caption={Implementación de Sensor Fusion en procesarLecturaSensor()}]
Future<void> procesarLecturaSensor(int ruido, double flujo) async {
  final DateTime timestampLectura = DateTime.now();

  // 1. Sensor Fusion
  final String estado;
  if (ruido > _umbralRuido && flujo > _umbralFlujoNormal) {
    estado = 'Uso Normal';
  } else if (ruido > _umbralRuido && flujo < _umbralFlujoFuga) {
    estado = 'Posible Fuga';
  } else {
    estado = 'Sin Clasificar';
  }

  _estadoActual = estado;
  caudalActual  = flujo;

  // 2. Buffer circular (maximo 10 lecturas)
  _bufferLecturas.add({
    'ruido': ruido, 'flujo': flujo,
    'estado': estado,
    'timestamp': timestampLectura.toIso8601String(),
  });
  if (_bufferLecturas.length > 10) _bufferLecturas.removeAt(0);

  // 3. Rate-limit para Groq
  if (estado == 'Posible Fuga') {
    _contadorPosibleFuga++;
  } else {
    _contadorPosibleFuga = 0;
  }

  if (_contadorPosibleFuga >= _umbralFugasConsecutivas && !iaProcesando) {
    final ahora = DateTime.now();
    final cooldownExpirado = _lastGroqEvaluation == null ||
        ahora.difference(_lastGroqEvaluation!).inMinutes >= _cooldownMinutos;

    if (cooldownExpirado) {
      _lastGroqEvaluation  = ahora;
      _contadorPosibleFuga = 0;
      await _llamarGroqPorFuga(flujo, timestampLectura);
    }
  }

  notifyListeners();
}
\end{lstlisting}

% -----------------------------------------------------------
\subsection{Mecanismo de Rate-Limiting para Groq}
% -----------------------------------------------------------

La API gratuita de Groq tiene límites de solicitudes por minuto. Para no saturarla,
el controlador implementa un doble mecanismo de protección:

\begin{enumerate}[leftmargin=2em]
  \item \textbf{Umbral de confirmación:} Solo se llama a Groq cuando el estado
        \texttt{"Posible Fuga"} se mantiene durante \textbf{3 lecturas consecutivas}
        (aproximadamente 15 segundos a una lectura cada 5 s). Una lectura aislada
        no desencadena una evaluación de IA.

  \item \textbf{Cooldown de 3 minutos:} Tras cada llamada exitosa a Groq, se registra
        el timestamp en \texttt{\_lastGroqEvaluation}. Ninguna llamada adicional será
        realizada hasta que hayan transcurrido \textbf{180 segundos}, incluso si el
        estado de fuga persiste.
\end{enumerate}

\begin{longtable}{p{4.5cm} p{9cm}}
\caption{Parámetros del mecanismo de rate-limiting --- La tabla define cada constante
  usada para controlar la frecuencia de llamadas a la API de Groq.}
\label{tab:rate-limit} \\
\toprule
\textbf{Constante} & \textbf{Valor y descripción} \\
\midrule
\endfirsthead
\toprule
\textbf{Constante} & \textbf{Valor y descripción} \\
\midrule
\endhead
\bottomrule
\endlastfoot
\texttt{\_umbralFugasConsecutivas} & \texttt{3} --- Número mínimo de lecturas
  consecutivas con ``Posible Fuga'' para activar Groq. \\[0.5em]
\texttt{\_cooldownMinutos} & \texttt{3} --- Minutos de espera mínima entre dos
  llamadas sucesivas a la API de Groq. \\
\end{longtable}

% ============================================================
\section{Guía de Integración para Proyectos Externos}
% ============================================================

Esta sección está dirigida a desarrolladores de proyectos externos (p.\,ej., sistemas
universitarios de análisis de datos) que deseen consumir el endpoint
\texttt{GET /api/history} de EdgeLeak~AI.

\subsection{Requisitos previos}

\begin{itemize}[leftmargin=2em]
  \item La aplicación EdgeLeak~AI debe estar \textbf{abierta y ejecutándose} en el
        teléfono (el servidor se detiene cuando la app se cierra).
  \item El teléfono y el equipo cliente deben estar \textbf{en la misma red WiFi local}.
  \item Conocer la \textbf{dirección IP local del teléfono} (ver cuadro de aviso en §3.2.4).
  \item Disponer de la \textbf{API Key} compartida por el equipo EdgeLeak~AI.
\end{itemize}

\subsection{Cambiar la API Key (recomendado en producción)}

Para mayor seguridad, el equipo de EdgeLeak~AI puede personalizar la API Key añadiendo
la variable \texttt{LOCAL\_API\_KEY} al archivo \texttt{.env} del proyecto Flutter:

\begin{lstlisting}[caption={Configurar API Key personalizada en .env}]
GROQ_API_KEY=gsk_xxxxxxxxxxxx
LOCAL_API_KEY=mi-clave-secreta-personalizada-2024
\end{lstlisting}

El proyecto cliente deberá usar el mismo valor en el header \texttt{x-api-key}.

\subsection{Ejemplo de integración en Python}

\begin{lstlisting}[language=Python, caption={Cliente Python para GET /api/history}]
import requests

# Configuracion
IP_TELEFONO = "192.168.1.100"  # Reemplazar con la IP real del telefono
API_KEY     = "edgeleak-share-2024"
URL         = f"http://{IP_TELEFONO}:8080/api/history"

def obtener_historial_fugas():
    headers = {"x-api-key": API_KEY}
    response = requests.get(URL, headers=headers, timeout=10)

    if response.status_code == 200:
        historial = response.json()
        print(f"Total de alertas: {len(historial)}")
        for alerta in historial:
            print(f"[{alerta['timestamp']}] "
                  f"{alerta['veredicto']} | "
                  f"Severidad: {alerta['severidad']} | "
                  f"{alerta['mensaje']}")
        return historial
    elif response.status_code == 401:
        print("Error 401: API Key incorrecta o ausente.")
    else:
        print(f"Error HTTP {response.status_code}: {response.text}")

if __name__ == "__main__":
    obtener_historial_fugas()
\end{lstlisting}

\subsection{Ejemplo de integración en JavaScript (Node.js)}

\begin{lstlisting}[language={}  , caption={Cliente Node.js para GET /api/history}]
const http = require("http");

const options = {
  hostname: "192.168.1.100",  // IP del telefono
  port: 8080,
  path: "/api/history",
  method: "GET",
  headers: {
    "x-api-key": "edgeleak-share-2024",
  },
};

http.get(options, (res) => {
  let data = "";
  res.on("data", (chunk) => { data += chunk; });
  res.on("end", () => {
    if (res.statusCode === 200) {
      const historial = JSON.parse(data);
      console.log(`Total alertas: ${historial.length}`);
      historial.forEach((a) => {
        console.log(`${a.timestamp} | ${a.veredicto} | ${a.severidad}`);
      });
    } else if (res.statusCode === 401) {
      console.error("Error 401: API Key invalida o ausente.");
    }
  });
});
\end{lstlisting}

\subsection{Ejemplo de integración en Dart/Flutter (otro proyecto)}

\begin{lstlisting}[language=Dart, caption={Cliente Dart para GET /api/history}]
import 'dart:convert';
import 'package:http/http.dart' as http;

Future<List<Map<String, dynamic>>> obtenerHistorialEdgeLeak({
  required String ipTelefono,
  required String apiKey,
}) async {
  final url = Uri.parse('http://$ipTelefono:8080/api/history');
  final response = await http.get(url, headers: {'x-api-key': apiKey});

  if (response.statusCode == 200) {
    final List<dynamic> lista = jsonDecode(response.body);
    return lista.cast<Map<String, dynamic>>();
  } else if (response.statusCode == 401) {
    throw Exception('API Key invalida o ausente.');
  } else {
    throw Exception('Error HTTP ${response.statusCode}');
  }
}

// Ejemplo de uso:
// final historial = await obtenerHistorialEdgeLeak(
//   ipTelefono: '192.168.1.100',
//   apiKey: 'edgeleak-share-2024',
// );
\end{lstlisting}

\subsection{Prueba rápida con curl}

Para verificar la conectividad desde cualquier terminal de línea de comandos:

\begin{lstlisting}[language={}, caption={Verificación de endpoints con curl}]
# Probar POST /api/sensor (simular envio de un ESP32)
curl -X POST http://192.168.1.100:8080/api/sensor \
  -H "Content-Type: application/json" \
  -d '{"ruido": 1800, "flujo": 0.04}'

# Probar GET /api/history con API Key correcta
curl http://192.168.1.100:8080/api/history \
  -H "x-api-key: edgeleak-share-2024"

# Probar GET /api/history sin API Key (debe devolver 401)
curl http://192.168.1.100:8080/api/history
\end{lstlisting}

% ============================================================
\section{Esquema de Base de Datos (Referencia)}
% ============================================================

El endpoint \texttt{GET /api/history} consume la tabla \texttt{historial} de la base
de datos SQLite local. A continuación se detalla el esquema completo de las tablas
relevantes del archivo \texttt{edgeleak\_v3.db}.

\subsection{Tabla: historial}

\begin{mdframed}[linecolor=normalgreen,linewidth=1pt,backgroundcolor=green!3]
\textbf{Descripción:} Almacena el registro cronológico de todas las alertas de fuga
generadas por el análisis de la IA (Groq). Cada fila representa una evaluación
completada, incluyendo el veredicto, la severidad y la recomendación del LLM. Es la
fuente principal de datos del endpoint \texttt{GET /api/history}.
\end{mdframed}

\vspace{0.5em}

\begin{longtable}{>{\ttfamily}p{3.5cm} p{2cm} p{2cm} p{5.5cm}}
\caption{Esquema de la tabla \texttt{historial} --- Registro cronológico de alertas de
  fuga generadas por la IA. Cada columna almacena un aspecto distinto de la evaluación.}
\label{tab:tabla-historial} \\
\toprule
\textbf{Columna} & \textbf{Tipo SQL} & \textbf{Nulo?} & \textbf{Descripción} \\
\midrule
\endfirsthead
\toprule
\textbf{Columna} & \textbf{Tipo SQL} & \textbf{Nulo?} & \textbf{Descripción} \\
\midrule
\endhead
\bottomrule
\endlastfoot
id & INTEGER & No &
  Clave primaria, autoincremental. Identifica de forma única cada alerta. \\[0.5em]
veredicto & TEXT & Sí &
  Clasificación de la IA. Valor esperado: \texttt{"Fuga Detectada"} o
  \texttt{"Flujo Normal"}. \\[0.5em]
severidad & TEXT & Sí &
  Nivel de urgencia asignado por el LLM: \texttt{"Crítica"}, \texttt{"Advertencia"}
  o \texttt{"Normal"}. \\[0.5em]
mensaje & TEXT & Sí &
  Texto en lenguaje natural con la recomendación de acción generada por el modelo
  Llama-3.3-70B de Groq. \\[0.5em]
fecha & TEXT & Sí &
  Fecha y hora de la detección almacenada como cadena ISO~8601
  (\texttt{YYYY-MM-DDTHH:mm:ss.sssZ}). \\
\end{longtable}

\subsection{Tabla: usuarios}

\begin{mdframed}[linecolor=primaryblue,linewidth=1pt,backgroundcolor=blue!3]
\textbf{Descripción:} Almacena las cuentas de usuario registradas en el sistema EdgeLeak AI.
Gestiona tanto los operadores como los administradores, con soporte para contraseñas
temporales que obligan al cambio en el primer inicio de sesión. Esta tabla es gestionada
exclusivamente por \texttt{DatabaseService} y no es expuesta por ningún endpoint HTTP.
\end{mdframed}

\vspace{0.5em}

\begin{longtable}{>{\ttfamily}p{3.5cm} p{2cm} p{2cm} p{5.5cm}}
\caption{Esquema de la tabla \texttt{usuarios} (versión 3) --- Registro de cuentas de
  usuario del sistema. Incluye los cuatro campos de nombre separados añadidos en la
  Actualización \#5 y el rol para control de acceso.}
\label{tab:tabla-usuarios} \\
\toprule
\textbf{Columna} & \textbf{Tipo SQL} & \textbf{Default} & \textbf{Descripción} \\
\midrule
\endfirsthead
\toprule
\textbf{Columna} & \textbf{Tipo SQL} & \textbf{Default} & \textbf{Descripción} \\
\midrule
\endhead
\bottomrule
\endlastfoot
id & INTEGER & --- &
  Clave primaria autoincremental. Identifica de forma única cada usuario. \\[0.5em]
primer\_nombre & TEXT & \texttt{''} &
  Primer nombre del usuario. Campo obligatorio en el formulario de registro. \\[0.5em]
segundo\_nombre & TEXT & \texttt{''} &
  Segundo nombre del usuario. Campo opcional. \\[0.5em]
primer\_apellido & TEXT & \texttt{''} &
  Primer apellido del usuario. Campo obligatorio en el formulario de registro. \\[0.5em]
segundo\_apellido & TEXT & \texttt{''} &
  Segundo apellido del usuario. Campo opcional. \\[0.5em]
correo & TEXT & --- &
  Dirección de correo electrónico. Restricción \texttt{UNIQUE}: no puede repetirse
  entre usuarios. Usado como identificador de inicio de sesión. \\[0.5em]
password & TEXT & --- &
  Contraseña del usuario almacenada como texto (MVP). En producción se recomienda
  almacenar el hash \texttt{bcrypt}. \\[0.5em]
es\_temporal & INTEGER & \texttt{0} &
  Indicador booleano (\texttt{0} = false, \texttt{1} = true). Si es \texttt{1},
  obliga al usuario a cambiar la contraseña en su primer inicio de sesión. \\[0.5em]
rol & TEXT & \texttt{'operador'} &
  Rol de acceso del usuario. Valores posibles: \texttt{'operador'} (solo lectura
  de alertas) o \texttt{'admin'} (acceso al CRUD de administración). \\
\end{longtable}

% ============================================================
\section{Consideraciones de Seguridad}
% ============================================================

\begin{itemize}[leftmargin=2em]
  \item \textbf{API Key como primera línea de defensa:} El endpoint \texttt{GET /api/history}
        está protegido mediante el header \texttt{x-api-key}. Aunque una API Key en
        texto claro no es el mecanismo más robusto posible, es suficiente para un
        contexto universitario en red local donde el riesgo de interceptación es bajo.
        En una fase productiva se recomienda migrar a OAuth~2.0 o JWT.

  \item \textbf{Ausencia de HTTPS:} El servidor HTTP local no implementa TLS/SSL. Las
        comunicaciones son en texto claro. Esto es aceptable en una red doméstica
        controlada, pero no debe exponerse a internet sin un proxy inverso con HTTPS
        (p.\,ej., nginx + Let's Encrypt).

  \item \textbf{Endpoint POST sin autenticación:} El endpoint \texttt{POST /api/sensor}
        no requiere autenticación. Esto asume que solo dispositivos de confianza
        (el ESP32 del proyecto) tienen acceso a la red local. En entornos con múltiples
        usuarios de red, se recomienda añadir un header \texttt{x-device-key} o filtrado
        por dirección IP.

  \item \textbf{Configuración de la API Key en .env:} El archivo \texttt{.env} no debe
        ser comprometido al repositorio de código. Asegurarse de que
        \texttt{.env} esté incluido en el archivo \texttt{.gitignore} del proyecto.

  \item \textbf{Puerto abierto en el dispositivo:} El servidor escucha en el puerto
        8080 de todas las interfaces de red del teléfono. En Android 10+, las apps
        pueden abrir puertos de servidor locales sin permisos especiales gracias al
        contexto de red local. Sin embargo, si el teléfono está conectado a una red
        corporativa, la política de red puede bloquear el puerto.
\end{itemize}

% ============================================================
\section{Resumen de Archivos Modificados}
% ============================================================

\begin{longtable}{p{7.5cm} p{6cm}}
\caption{Resumen de archivos modificados en la Actualización \#6 --- La tabla lista
  cada archivo afectado y describe el cambio principal introducido en esta actualización.}
\label{tab:archivos-resumen} \\
\toprule
\textbf{Archivo} & \textbf{Cambio principal} \\
\midrule
\endfirsthead
\toprule
\textbf{Archivo} & \textbf{Cambio principal} \\
\midrule
\endhead
\bottomrule
\endlastfoot
\texttt{lib/data/services/local\_api\_service.dart} &
  \textbf{Nuevo.} Servidor HTTP local con \texttt{POST /api/sensor} y
  \texttt{GET /api/history}. \\[0.5em]
\texttt{lib/controllers/dashboard\_controller.dart} &
  \textbf{Modificado.} Sensor Fusion, buffer circular, rate-limit de Groq,
  instanciación y ciclo de vida de \texttt{LocalApiService}. \\
\end{longtable}

\end{document}
